\section{Conclusions}
\label{sec:conclusions}

This project set out to determine whether wrist-worn actigraphy carries sufficient signal to distinguish bipolar from unipolar depression using deep sequence models, without relying on hand-crafted feature engineering. Through a systematic four-approach investigation, we find a clear answer: during depressive episodes, the motor activity patterns of bipolar and unipolar patients in the Depresjon dataset are statistically indistinguishable. Our most rigorous classifier achieves a best honest accuracy of 63.4\% using 48-hour CNN-LSTM windows, barely above the 65.2\% majority-class baseline, and direct hypothesis testing confirms negligible group differences ($p = 0.776$, Cohen's $d = {-}0.060$ \cite{cohen1988statistical}). This is a rigorous negative result. It does not indicate that actigraphy is uninformative about mental health more broadly; it specifically rules out wrist activity variability during depressive episodes as a standalone discriminator between bipolar and unipolar depression. The failure of a 223,682-parameter deep network \cite{paszke2019pytorch} to outperform a 5-feature logistic regression \cite{pedregosa2011sklearn} is itself an instructive finding: in clinical machine learning, more complex models do not compensate for insufficient or non-discriminative data. The negative result redirects future research toward concrete and well-motivated alternatives.

Three directions stand out as most promising. First, the Depresjon dataset captures only depressive episodes, but the defining clinical distinction of bipolar disorder is the existence of manic or hypomanic phases \cite{goodwin2007manic}. Longitudinal monitoring across full mood episode cycles, including periods of elevated mood, reduced sleep need, and increased goal-directed activity, would provide the class-discriminating signal that depressive-phase recordings cannot. Wehr et al.\ \cite{wehr1987sleep} document that bipolar patients exhibit phase-advanced circadian rhythms most visibly during manic states, not depressive ones, which is precisely the signal that is absent from this dataset. Second, wrist actigraphy alone is a narrow window into the physiology underlying mood disorders. Integrating additional biosignal modalities would substantially enrich the available information: polysomnographic sleep stage data can capture REM-sleep abnormalities that differ between bipolar and unipolar depression at a mechanistic level; heart rate variability reflects autonomic nervous system dysregulation that tracks mood state; and ecological momentary assessment, which captures self-reported mood several times daily, provides a subjective signal that can be aligned with objective sensor data to identify the lag between physiological change and subjective experience. Third, and most fundamentally, any future study targeting the bipolar-versus-unipolar discrimination task with statistical credibility requires a substantially larger bipolar cohort. To achieve 80\% power for detecting a small effect size ($d = 0.2$) at the $\alpha = 0.05$ significance level, approximately 200 bipolar participants are needed \cite{cohen1988statistical}. Such a cohort must also reflect demographic diversity in age, sex, medication history, and occupation to ensure that learned features reflect genuine diagnostic signal rather than confounded population differences.

The ethical dimensions of this line of research deserve careful consideration alongside the technical findings. Diagnostic algorithms trained on small, demographically constrained datasets carry the risk of encoding spurious patterns that fail to generalize to broader patient populations, potentially automating misdiagnosis at scale. The Depresjon dataset, with only 8 bipolar participants, is not large enough to characterize the true distribution of bipolar motor activity patterns across the condition's full heterogeneity, and any deployed system trained on it could exhibit systematic failures for subgroups not adequately represented in the sample. Beyond generalization, continuous passive wrist monitoring raises meaningful privacy concerns: the data collected is detailed enough to infer not just mood state but daily routines, social patterns, and behavioral habits, and patients in vulnerable populations may not fully appreciate the scope of what is recorded. Algorithmic accountability is a further concern: clinicians may implicitly defer to model outputs rather than exercising independent judgment, making interpretability and calibration of such systems safety-critical. Our best models reported 88\% accuracy while detecting zero bipolar cases, a failure mode that would not be obvious from the headline metric alone. Any future deployment of actigraphy-based classifiers must be accompanied by transparent reporting of sensitivity, specificity, calibration, and explicit uncertainty quantification alongside headline accuracy figures. The path from a promising result on a 55-person dataset to a clinically validated tool is long and uncertain, and intellectual honesty about negative findings such as ours is an essential part of navigating it responsibly.
